\begin{table}[!h]
\centering
\caption{\label{tab:varselect7}Selección de rezagos del VAR de 7 variables}
\centering
\begin{threeparttable}
\begin{tabular}[t]{ccccc}
\toprule
Rezagos & AIC & HQ & BIC & Portmanteau (18)\\
\midrule
1 & \textbf{-2.817} & \textbf{-2.530} & \textbf{-2.101} & 0.009\\
2 & -2.814 & -2.276 & -1.471 & 0.043\\
3 & -2.676 & -1.887 & -0.708 & 0.035\\
4 & -2.571 & -1.531 & 0.024 & 0.015\\
5 & -2.457 & -1.166 & 0.765 & 0.029\\
\addlinespace
6 & -2.351 & -0.808 & 1.497 & 0.012\\
7 & -2.237 & -0.444 & 2.237 & 0.016\\
8 & -2.089 & -0.045 & 3.011 & 0.003\\
\bottomrule
\end{tabular}
\begin{tablenotes}
\item \textit{Note: } 
\item Menor valor = rezago óptimo (en negrita). BIC = criterio de Schwarz, HQ = Hannan-Quinn. La última columna reporta el $p$-valor del test de Portmanteau ajustado a 18 rezagos (H0: ausencia de autocorrelación). Los criterios seleccionan un rezago, pero éste deja residuos autocorrelacionados; se adopta $p=2$, la especificación más parsimoniosa que mitiga la autocorrelación.
\end{tablenotes}
\end{threeparttable}
\end{table}